%% kapa: exam=1 round=23 year=2012 period=1 form=A subject=economics subject_ko=경제원론 seq=24 number=64
%% source: Q-Net 감정평가사 기출문제 (cst00309) artlSeq=5052578 file=raw/1008/5052578/2012년 제23회 감정평가사 제1차시험 기출문제(1교시 A형).pdf
%% conversion: kordoc 4.24.51 -> tools/gen1.py (auto); review=proofread-sonnet; flags=none
\begin{question}{64}
폐쇄경제 IS-LM 모형과 관련된 설명으로 옳은 것은? (단, IS 곡선은 우하향, LM 곡선은 우상향한다.)
\begin{choices}
  \item IS 곡선과 LM 곡선에서 총공급 곡선이 도출된다.
  \item 정부지출의 구축효과는 발생하지 않는다.
  \item 현재 경제상태가 IS 곡선의 왼쪽, LM 곡선의 오른쪽에 있다면 상품시장은 초과공급, 화폐시장은 초과수요 상태이다.
  \item 피구효과(Pigou effect)에 의하면 물가수준이 하락할 때 IS 곡선이 우측으로 이동하여 국민소득이 증가한다.
  \item IS-LM 모형에서 물가수준은 내생변수이다.
\end{choices}
\end{question}
\answer{4}
